\section{Introduction} \label{sec:41-research02-introduction}

\begin{foldable}
  The dual challenge of domain-coverage deficit and information-bottleneck deficit is qualitatively different from the few-shot setting addressed in Chapter~\ref{ch:30-research01}.
  Few-shot black-box \ac{KD} at least provides a distributional anchor: real training images from the teacher's domain, even if limited in diversity.
  Data-free black-box \ac{KD} removes this anchor entirely.
  The generator must reconstruct the image manifold from hard-label queries alone, with no real samples to guide synthesis and no soft targets to preserve inter-class structure.
\end{foldable}

\begin{foldable} % Recent methods
  Recent methods in BBDFKD~\cite{wang2021zero,zhang2022ideal,yuan2024data} use synthetic data, but they still encounter two primary bottlenecks.
  First, they suffer from a lack of domain coverage: synthetic samples fail to replicate the hierarchical structures and semantic variation of natural environments, resulting in a distribution mismatch with the teacher's expertise.
  Second, the distillation signals are suboptimal.
  By relying exclusively on hard top-1 indices, students are denied the \textit{dark} knowledge of inter-class relationships---a cornerstone of \ac{KD} essential for emulating expert logic.
\end{foldable}

\begin{foldable} % Our method
  \textbf{Our method.} We propose \textbf{\acl{DIPKD}} (\acs{DIPKD}) to tackle the BBDFKD problem, \ie, training a student from a black-box teacher with only top-1 signal and without access to any real data.
  To improve \ac{KD}, we incorporate \textit{diversity}, \ie, how broadly the samples distribute, into a collaborative three-phase pipeline:{
    (1)~\textit{Synthesis}: crafts synthetic \textit{image priors} with diverse patterns and semantics of natural objects;
    (2)~\textit{Contrast}: employs a \textit{primer student} to act as a white-box mediator and ensure synthetic samples are informationally distinct via contrastive learning;
    and (3)~\textit{Distillation}: distills the final student optimally with hard labels from teacher and soft probabilistic labels from the primer student.
  }
  Synthesis and Contrast together address the domain-coverage deficit; the primer student's soft labels in Distillation partially restore the information-bottleneck deficit.
\end{foldable}

\begin{foldable} % Contribution
  \textbf{Contribution.} In summary, we highlight our contributions as follows:
  \begin{enumerate}
    \item We introduce \textit{image priors}, a novel class of synthetic images exhibiting diversity across hierarchical structures, nonlinearity, and semantics.
    \item We propose the novel use of a \textit{primer student}, enabling contrastive optimization and extraction of soft probabilistic signals originally restricted in black-box \ac{KD}.
    \item Extensive analysis of our framework demonstrates that it effectively expands the semantic coverage of synthetic domains, leading to our state-of-the-art performance across 12 benchmarks.
  \end{enumerate}
\end{foldable}

\begin{foldable}
  The remainder of this chapter develops the framework (\S\ref{sec:43-research02-framework}), reports extensive experiments and ablations (\S\ref{sec:44-research02-experiments}), and concludes with discussion (\S\ref{sec:45-research02-conclusion}).
\end{foldable}
